% =====================================================================================
%  Chapter 4 — Control Architectures
%  Thesis: Comparison of Control Strategies for Interaction-Force Aware Multirotor Teams
%
%  DRAFT, 23 August 2026 (writing-track item W.8).
%  REVISED 2026-10-14 (thesis-revision run, round 1).
%  REVISED 2026-10-15 (round 2: final study controllers). Rewritten around ONE controller with
%  three switches (Strategies 0-3) and a separate, not implemented Strategy 4 (FBL). Markers:
%  % NEW, % CHANGED, % DECISION, % VERIFY, % CONFIRM-AFTER-LAB. Change log:
%  docs/thesis_revision/round2/round2_change_log.md.
%
%  THE RULES THIS CHAPTER MUST NOT BREAK.
%  (1) The visible text names no implementation, person, language, file or code identifier; it
%      describes the controllers by function and cites the papers whose equations they follow.
%      Code locations appear only in comments, for traceability.
%  (2) Strategy 4 is not implemented and is never described as if it existed.
%  (3) No result of any kind (no flown, lab, replay or simulation number).
%
%  Code basis (comments only): crazyflie-firmware/src/modules/src/controller/controller_omar_indi.c
%  (+ .h), src/modules/src/ns2_residual.c (+ .h), src/platform/interface/platform_defaults_cf21bl.h;
%  firmware_app/traj_iface.c (rnn parameter group), firmware_app/app-config-bl (CF_MASS);
%  crazyswarm2/crazyflie/config/study_strategies.json and study/tree_S0..S3.patch.
%
%  Label note: sec:ctrl-s3 now labels the hybrid (Strategy 3); the feedback-linearising strategy,
%  formerly Strategy 3, is Strategy 4 and carries the new label sec:ctrl-s4.
% =====================================================================================

\chapter{Control Architectures}
\label{ch:control}

Chapter~\ref{ch:model} defined the residual wrench and the two ways of obtaining it, by measurement
and by prediction. This chapter describes the control strategies compared in this thesis and, for
each, the point at which knowledge of the interaction enters the control law.

% CHANGED (round 2): the strategy list of the final study.
Five strategies are defined (Table~\ref{tab:ctrl-summary}). Four of them, Strategies~0--3, are one
controller: a geometric tracking controller to which three independent switches add a measured
residual force, a measured residual moment and a learned residual force. Strategy~0 has all three
switches off and is the uncompensated baseline (``uncompensated'' means that it does not use the
residual); Strategy~1 uses the measured residual (incremental nonlinear dynamic inversion, INDI);
Strategy~2 uses the learned residual; Strategy~3 uses both. Strategy~4 is a different controller, a
feedback-linearising law with a learned residual, and is described as a reference design.
% CHANGED: code, controller_omar_indi.c (indi bit 0 / bit 1, L231/L365 region; rnn.en via
%   ns2_active() L241-249); study_strategies.json blocks S0-S3 (indi 0/3 x rnn.en 0/1).

% DECISION D-05 (H9): the last column states implementation in the flight code, not flight status;
%   nothing has been flown with the final controllers. Options: keep this neutral column, or delete it.
\begin{table}[t]
  \centering
  \caption{The strategies. Switches A and B add the measured residual force and moment, switch C
    the learned residual force; NN denotes the learned residual and FBL feedback linearisation. ``Knows'' refers to the residual: whether the controller has a
    measurement of it, a prediction of it, both, or neither. ``Enters at'' is the quantity the
    residual is subtracted from, and ``Rotor speed'' says whether the strategy needs rotor-speed
    sensing at run time.}
  \label{tab:ctrl-summary}
  \small
  \begin{tabular}{@{}cllccc@{}}
    \toprule
    \# & Strategy & Knows & Switches & Rotor speed & Implemented \\
    \midrule
    0 & Geometric baseline           & nothing             & ---     & no                      & yes \\
    1 & INDI                         & measured            & A, B    & yes                     & yes \\
    2 & Geometric + NN               & predicted           & C       & no\textsuperscript{a}   & yes \\
    3 & INDI + NN                    & measured, predicted & A, B, C & yes                     & yes \\
    4 & FBL + NN                     & predicted           & ---     & no\textsuperscript{a}   & no \\
    \bottomrule
  \end{tabular}

  \vspace{0.4em}
  \begin{minipage}{0.9\linewidth}\footnotesize
    \textsuperscript{a}~Not at run time. Producing the training labels requires the measured
    residual \eqref{eq:model-ares}, and therefore rotor-speed sensing on the vehicles that collect
    the data, as for the learned method of Cobo-Briesewitz et al.~\cite{cobobriesewitz2026lindi}.
    The evaluation of Chapter~\ref{ch:setup} records rotor speeds on every flight.
  \end{minipage}
\end{table}

% CHANGED (round 2): the fairness argument reverses (author decision 6).
Because Strategies~0--3 are one controller, they share the position and attitude gains, the
filters and the physical parameters, and differ only in which switches are set. What makes the
comparison fair is therefore both the experimental conditions held constant and the shared control
law; Section~\ref{sec:ctrl-fairness} states precisely what is shared and what the comparison still
does not guarantee.

Each strategy is described in two parts, kept separate: the law that is flown, and the published
method it follows, with the deviations from that method listed.

% -------------------------------------------------------------------------------------
\section{The common controller}
\label{sec:ctrl-common}

% CHANGED (round 2): one controller for Strategies 0-3.
Strategies~0--3 share one position loop and one attitude loop. The switches add terms to them; they
do not replace either loop. The equations are stated once in this section.

\subsection{Position loop}

The desired acceleration is
%
\begin{equation}
  a_{d}
  \;=\; a_{\mathrm{ff}} + g e_3 + K_p e_p + K_v e_v + k_{i,z}\, e_3 \!\int\! e_3^\top e_p \,\mathrm{d}t ,
  \label{eq:ctrl-ad}
\end{equation}
%
with $e_p = p_{\mathrm{des}} - p$, $e_v = v_{\mathrm{des}} - v$ and $a_{\mathrm{ff}}$ the
trajectory's acceleration. The integral acts on the vertical axis only, and its state is limited.
The switches then form the commanded acceleration
%
\begin{equation}
  a_{\mathrm{cmd}}
  \;=\; a_d
      \;\underbrace{-\; a_{\mathrm{res},f}}_{\text{switch A}}
      \;\underbrace{-\; \hat{a}_{\mathrm{res}}}_{\text{switch C}},
  \label{eq:ctrl-ades}
\end{equation}
%
where $a_{\mathrm{res},f}$ is the measured residual \eqref{eq:model-ares} after the filters of
Section~\ref{sec:ctrl-s1} and $\hat a_{\mathrm{res}}$ the learned prediction of
Section~\ref{sec:ctrl-s2}. Both residual terms are subtracted, the sign derived in
Section~\ref{sec:model-sign}. The part of \eqref{eq:ctrl-ad} without $g e_3$ is the $a_{\mathrm{des}}$
of Chapter~\ref{ch:model}.
% CHANGED (round 2; D-01 and D-02 resolved): code, controller_omar_indi.c L214-249 (acc_d with
%   +GRAVITY_MAGNITUDE, Kpos_D*vel_e + Kpos_P*pos_e + i_term; z-only integral when Kpos_I.z != 0,
%   state clamped to +-Kpos_I_limit = 2; NS2 a_d - hold, L241-249), L250-268 (a_indi = a_rpm_f - a_imu_f
%   = -a_res_f, thrust = m (a_d + a_indi) . R e3). The NS2 term is subtracted before the INDI term is
%   added; the sum is the same. The integral states are reset while the commanded thrust is below
%   0.01 N (on the ground, L271-274) and the vertical one also outside position mode (L201-206).
%   No velocity-error conditioning exists in this controller (round-1 text removed).

The commanded thrust vector is $m\,a_{\mathrm{cmd}}$. Its projection on the body $z$-axis is the
collective thrust, $f = m\,a_{\mathrm{cmd}}^\top R e_3$, and its direction, together with the desired
yaw $\psi_{\mathrm{des}}$, fixes the desired attitude: with
$y_c = [-\sin\psi_{\mathrm{des}},\, \cos\psi_{\mathrm{des}},\, 0]^\top$, the columns of
$R_{\mathrm{des}}$ are
\begin{equation*}
  x_{\mathrm{des}} = \frac{y_c \times a_{\mathrm{cmd}}}{\lVert y_c \times a_{\mathrm{cmd}} \rVert}, \qquad
  y_{\mathrm{des}} = \frac{a_{\mathrm{cmd}} \times x_{\mathrm{des}}}{\lVert a_{\mathrm{cmd}} \times x_{\mathrm{des}} \rVert}, \qquad
  z_{\mathrm{des}} = x_{\mathrm{des}} \times y_{\mathrm{des}},
\end{equation*}
so that the body $z$-axis points along $a_{\mathrm{cmd}}$ and the desired heading fixes the rotation
about it. Because the residual terms are inside $a_{\mathrm{cmd}}$, they change both the thrust and
the desired attitude.
% CHANGED: code, controller_omar_indi.c L276-282 (xb = normalize(yc x F_d), yb = normalize(F_d x xb),
%   zb = xb x yb).

\subsection{Attitude loop}

% CHANGED (round 2): the full geometric law of Lee et al. with an attitude integral; the switch-B
%   correction is added to it (Section ctrl-s1).
The attitude error and the rate error are those of Lee et al.,
$e_R = \tfrac{1}{2}\left(R_{\mathrm{des}}^\top R - R^\top R_{\mathrm{des}}\right)^\vee$ and
$e_\omega = \omega - \omega_r$ with $\omega_r = R^\top R_{\mathrm{des}}\,\omega_{\mathrm{des}}$, where
the desired angular velocity $\omega_{\mathrm{des}}$ and angular acceleration
$\dot\omega_{\mathrm{des}}$ follow from the trajectory's jerk and snap by differential flatness. The
torque is
%
\begin{equation}
  \tau_{\mathrm{geo}} = -K_R e_R - K_\omega e_\omega - K_I \!\int\! e_R \,\mathrm{d}t + \omega \times J\omega
         - J\!\left( \hat{\omega}\, \omega_r - R^\top R_{\mathrm{des}}\, \dot\omega_{\mathrm{des}} \right),
  \label{eq:ctrl-tau-geo}
\end{equation}
%
and switch~B adds to it the measured-moment correction of Section~\ref{sec:ctrl-s1}. Here $K_R$ and
$K_\omega$ are diagonal gains from attitude and rate error to torque, in \si{\newton\metre} per unit
of $e_R$ and \si{\newton\metre\second} per unit of $e_\omega$, and $K_I$ is the gain of an integral
of the attitude error. Apart from that integral, \eqref{eq:ctrl-tau-geo} is the torque law of the
source with its feedforward bracket, given in Section~\ref{sec:ctrl-s0}.
% CHANGED: code, controller_omar_indi.c L318-364 (eR from eRM; omega_des from jerk via B1, B3, C3,
%   D1-D3; omega_des_dot from snap via E1-E3; omega_r = R^T R_des omega_des; u = -KR eR - Komega ew
%   - KI i_error_att + omega x J omega - J (omega^ omega_r - R^T R_des omega_des_dot)). The attitude
%   integral accumulates without a limit and is reset on the ground (controllerOmarIndiReset).

\subsection{Parameters}

% NEW (round 2): the shared parameter set, from the code (author decision 6).
Table~\ref{tab:ctrl-params} lists the parameters of the common controller. All of them are the same
in Strategies~0--3.
\begin{table}[t]
  \centering
  \caption{Parameters of the common controller, identical in Strategies~0--3.}
  \label{tab:ctrl-params}
  \small
  \begin{tabular}{@{}lll@{}}
    \toprule
    Quantity & Symbol & Value \\
    \midrule
    Position gain, all axes & $K_p$ & \SI{7}{\per\second\squared} \\
    Velocity gain, all axes & $K_v$ & \SI{4}{\per\second} \\
    Vertical integral gain & $k_{i,z}$ & $1.5$\,s$^{-3}$ \\
    Attitude gain & $K_R$ & $\mathrm{diag}(0.007,\,0.007,\,0.008)$\,\si{\newton\metre} \\
    Rate gain & $K_\omega$ & $\mathrm{diag}(0.00115,\,0.00115,\,0.002)$\,\si{\newton\metre\second} \\
    Attitude integral gain & $K_I$ & $0.03$ on each axis \\
    Mass & $m$ & \SI{42.7}{\gram} \\
    Inertia & $J$ & $\mathrm{diag}(16.57,\,16.66,\,29.26)\times 10^{-6}$\,\si{\kilo\gram\metre\squared} \\
    Rotor thrust constant & $k_t$ & $3.912\times 10^{-8}$\,\si{\newton\second\squared} per rotor \\
    Arm length (centre to rotor) & $l$ & \SI{0.050}{\metre} \\
    Rotor drag torque to thrust & $\kappa$ & \SI{0.00569}{\metre} \\
    Residual filters & --- & second-order Butterworth, \SI{30}{\hertz} \\
    \bottomrule
  \end{tabular}
\end{table}
% CHANGED: code. Gains: controller_omar_indi.c L56-73 (Kpos_P 7, Kpos_D 4, KR, Komega, KI 0.03);
%   Kpos_Iz 1.5 from study_strategies.json S0-S3 / study/tree_S*.patch (default 0). Mass: CF_MASS =
%   CONFIG_MODIFIED_CF_MASS 42700 (firmware_app/app-config-bl:48-49). Inertia: L59-62, comment
%   "System Identification of the Crazyflie 2.0 Nano Quadrocopter, BA thesis, J. Foerster, ETHZ" --
%   i.e. the identification of a DIFFERENT airframe (Crazyflie 2.0), see the deviation in Sec. ctrl-s0
%   and round2_open_items.md. k_t: MOTORRPM2FORCE 3.91194027330774e-08 N/(rad/s)^2; arm: ARM_LENGTH
%   0.050 m; kappa: THRUST2TORQUE 0.00569278844371417 m (platform_defaults_cf21bl.h:45, 65-67).
%   Filters: L131-136 (cutoff 30 Hz, sample time 1/ATTITUDE_RATE = 2 ms). Units of k_{i,z}: the
%   integral of a position error times 1/s^3 gives an acceleration.

\subsection{Rates}

% CHANGED (round 2): rates of the final controller.
The controller executes at \SI{500}{\hertz}, on every second tick of the \SI{1}{\kilo\hertz} main
loop of the flight software. The motion-capture pose reaches the vehicles at \SI{100}{\hertz}, and
the learned residual is evaluated on every tenth main-loop tick, that is \SI{100}{\hertz}, and held
in between.
% CHANGED: code, controller_omar_indi.c L156 (RATE_DO_EXECUTE(ATTITUDE_RATE = 500 Hz)); stabilizer
%   main loop RATE_MAIN_LOOP 1000 Hz (stabilizer_types.h:383); ns2_residual.c ns2_update (tick % div,
%   rnn.div default 10, main-loop tick).

\subsection{Runtime selection}
\label{sec:ctrl-runtime}

% CHANGED (round 2): three switches; one parameter set per strategy.
The switches are parameters set at run time, not compile-time options: switch~A adds the measured
residual force in \eqref{eq:ctrl-ades}, switch~B the measured-moment correction to
\eqref{eq:ctrl-tau-geo}, and switch~C the learned residual in \eqref{eq:ctrl-ades}. Strategy~0 has
none set, Strategy~1 has A and B, Strategy~2 has C, and Strategy~3 has A, B and C. Each strategy
is one frozen set of parameter values, sent to the vehicle before take-off and read back from it.
The learned prediction is computed and recorded whatever switch~C says; the switch decides only
whether it enters \eqref{eq:ctrl-ades}. Switches A and B need the rotor-speed source and are inert
without it.
% CHANGED: code, controller_omar_indi.c (indi & 1, indi & 2, both gated by rpm_deck_available =
%   deck.bcRpm == 1, L179/L231/L365); ns2_residual.c ns2_update (always evaluates and logs) and
%   ns2_active (rnn.en); study_strategies.json (frozen blocks, read-back check before take-off).

The two vehicles of a formation run the same controller. The upper vehicle always flies
Strategy~0; the strategy under comparison, including any learned term, is flown by the lower
vehicle, the one that meets the downwash.
% CHANGED: code, study/tree_S0..S3.patch (top vehicle: indi 0, rnn.en 0 in every strategy).

Setting A or B alone gives two further configurations. They are diagnostic, used in the two-robot
results chapter to attribute an effect to the force or the moment term, and are not strategies.

% CONFIRM-AFTER-LAB (a): take-off and landing of the INDI strategies. A take-off with switches A and B
%   set from the ground has not been flown; the fallback is to take off with the geometric law and
%   set the switches at altitude. The yaml (tree_S1/S3.patch) currently pins the controller with
%   indi 3 for the whole flight, including take-off and landing.
How take-off and landing are flown for the strategies that use switches A and B is fixed at the
hardware gate (Section~\ref{sec:setup-gate}); where it differs from the flight phase compared, the
analysis windows of Section~\ref{sec:setup-metrics} exclude it.

Runtime selection matters more than it may appear. Comparing strategies that require different
firmware introduces a confound that cannot be ruled out afterwards: a flash cycle changes more
than the control law, and strategies flown on different days meet different battery states,
temperatures and motion-capture calibrations. With one controller and run-time switches the
strategies can be interleaved within a session, which makes the comparison a paired one.

% -------------------------------------------------------------------------------------
\section{Strategy 0: uncompensated geometric baseline}
\label{sec:ctrl-s0}

\paragraph{The law.} Equations~\eqref{eq:ctrl-ad}--\eqref{eq:ctrl-tau-geo} with all three switches
off. The controller has no knowledge of the interaction.

\paragraph{Reference.} Geometric tracking on $SE(3)$~\cite{lee2010geometrictracking}, which tracks
a position trajectory together with a single heading direction, four outputs, not a full
six-degree-of-freedom pose. The attitude error of \eqref{eq:ctrl-tau-geo} is that work's
$e_R = \tfrac{1}{2}(R_d^\top R - R^\top R_d)^\vee$.
Lee et al.\ write the translational dynamics with the third axis pointing \emph{down},
$m\dot v = mge_3 - fRe_3$, and define the position error as $e_x = x - x_d$, the opposite of the
$e_p = p_{\mathrm{des}} - p$ used here~\cite{lee2010geometricse3}. Position-loop terms therefore change
sign when transcribed into this thesis' east-north-up convention, whereas the attitude error and
the torque law, written in body coordinates, carry over unchanged.
The torque law of the source is
%
\begin{equation}
  M = -k_R e_R - k_\Omega e_\Omega + \Omega \times J\Omega
      - J\!\left( \hat{\Omega}\, R^\top R_d \Omega_d - R^\top R_d \dot{\Omega}_d \right),
  \label{eq:ctrl-tau-lee}
\end{equation}
%
with $e_\Omega = \Omega - R^\top R_d \Omega_d$; the symbols map to those used here as $k_R \to K_R$,
$k_\Omega \to K_\omega$, $\Omega \to \omega$, $R_d \to R_{\mathrm{des}}$ and $\Omega_d \to \omega_{\mathrm{des}}$.

% CHANGED (round 2): deviations of the flown law from the source.
\paragraph{Deviations.} Equation~\eqref{eq:ctrl-tau-geo} contains every term of
\eqref{eq:ctrl-tau-lee} and adds an integral of the attitude error; the position loop adds a
vertical integral; the desired attitude is built from the direction of $a_{\mathrm{cmd}}$ and the
desired heading as above, rather than by the heading projection of the source; and $\omega_{\mathrm{des}}$
and $\dot\omega_{\mathrm{des}}$ are computed from the trajectory's jerk and snap. The integral terms
mean that no stability result of the source transfers to the flown law unexamined. The inertia
used by the controller is the one identified for the earlier Crazyflie~2.0 airframe, not one
identified for the vehicles flown here; an inertia error affects the feedforward and gyroscopic
terms of \eqref{eq:ctrl-tau-geo} and, with switch~B, the moment correction.
% CHANGED: code, controller_omar_indi.c L59-62 (comment: Crazyflie 2.0 identification).
% VERIFY: the inertia source (Foerster, BA thesis, ETH Zurich) is not held; see
%   round2_proposed_references.md. Whether a Crazyflie 2.1 Brushless identification should replace
%   it is an author decision (round2_open_items.md).

\paragraph{What the source proves.} Not global asymptotic tracking. With the attitude error
function $\Psi$, exponential stability is established on $\Psi(0) < 1$, initial attitude errors
below \SI{90}{\degree}, and almost-global exponential attractiveness on $1 \le \Psi < 2$,
below \SI{180}{\degree}. This thesis claims neither for the flown law: the baseline is flown as a
reference condition, not as a certified controller.
Here $\Psi(R, R_d) = \tfrac{1}{2}\operatorname{tr}(I - R_d^\top R)$, which equals $1 - \cos\theta$
for a rotation by the angle $\theta$ between $R$ and $R_d$, so $\Psi = 1$ at \SI{90}{\degree} and
$\Psi = 2$ at \SI{180}{\degree}~\cite{lee2010geometricse3}.
% VERIFY: the regions are attributed to the CDC paper (no PDF held); they are checked against the
%   held companion arXiv:1003.2005v4 (Propositions 3 and 4 there).

% CHANGED (round 2; D-04 resolved): the role, and why training data come from dedicated flights.
\paragraph{Role.} Strategy~0 is the uncompensated baseline of Table~\ref{tab:ctrl-summary}. It is
flown in every scenario so that ``how much did compensation help'' has a denominator, and so that
the interaction can be separated from the model error discussed in
Section~\ref{sec:model-measure-fail}. It is also the controller of the dedicated training flights
from which the learned residual is obtained (Section~\ref{sec:ctrl-s2}); those flights are separate
from the comparison flights, so that the learned model is never evaluated on its own training data.

% -------------------------------------------------------------------------------------
\section{Strategy 1: INDI}
\label{sec:ctrl-s1}

% CHANGED (round 2; D-11 and D-01 resolved): the measured residual force and moment are added to the
%   common controller; the round-1 incremental attitude law of the retired implementation is removed.
\paragraph{The law.} Switches A and B. Switch~A subtracts the measured residual force in
\eqref{eq:ctrl-ades}. It is formed as in \eqref{eq:model-ares} from two accelerations: the
acceleration the rotors are producing, computed from the four measured rotor speeds,
\begin{equation*}
  a_{\mathrm{rot}} = \frac{1}{m}\Big(\sum_{k=1}^{4} k_t \varpi_k^2\Big) R e_3 - g e_3 ,
\end{equation*}
and the measured acceleration $a_{\mathrm{meas}}$, the accelerometer reading rotated into the world
frame with the estimated attitude and with gravity removed. Both pass identical second-order Butterworth low-pass filters at \SI{30}{\hertz},
sampled at \SI{500}{\hertz}, and $a_{\mathrm{res},f} = a_{\mathrm{meas},f} - a_{\mathrm{rot},f}$.
Switch~B adds to the torque \eqref{eq:ctrl-tau-geo} the measured-moment correction
%
\begin{equation}
  \tau = \tau_{\mathrm{geo}} + \big(\tau_{\mathrm{cur},f} - J \alpha_f\big),
  \label{eq:ctrl-s1-torque}
\end{equation}
%
where $\tau_{\mathrm{cur}}$ is the torque the rotors are producing, computed from the four rotor
thrusts $k_t\varpi_k^2$, the arm length and the rotor drag-to-thrust ratio, and $\alpha$ is the
angular acceleration obtained by differencing the gyro; both pass the same \SI{30}{\hertz} filter.
The correction is added to the full geometric torque; it does not replace the torque law.
% CHANGED: code, controller_omar_indi.c L179-196 (t_k = MOTORRPM2FORCE * (rpm_k * 2 pi / 60)^2 from the
%   deck log variables rpm.m1..m4, raw), L250-262 (a_rpm, a_imu = 9.81 * state->acc, filtered;
%   state->acc = R * accelerometer - g, kalman_core.c:790-797),
%   a_indi = a_rpm_f - a_imu_f), L365-410 (tau_rpm from t_k with arm 0.707 * ARM_LENGTH and
%   THRUST2TORQUE; angular_acc = (omega - omega_prev)/dt; tau_gyro = J angular_acc; both filtered;
%   u += tau_rpm_f - tau_gyro_f).

\paragraph{Why the terms reject the disturbance.} For the force, substituting
$a_{\mathrm{res},f}$ into $f R e_3 = m\,a_{\mathrm{cmd}}$ gives
\begin{equation*}
  f R e_3 = m\,(a_{\mathrm{des}} + g e_3) - m\,(a_{\mathrm{meas},f} - a_{\mathrm{rot},f})
          = f_{\mathrm{cur},f} R e_3 + m\,(a_{\mathrm{des}} - a_{\mathrm{meas},f}),
\end{equation*}
with $f_{\mathrm{cur}} = \sum_k k_t \varpi_k^2$ the thrust being produced: the new thrust vector is
the current one corrected by the acceleration error, the linear-acceleration increment of Tal and
Karaman and of Smeur et al.~\cite{tal2018indi,smeur2017cascadedindi}. For the moment, defining the
reference angular acceleration that the geometric law commands,
\begin{equation}
  J\alpha_{\mathrm{ref}} = \tau_{\mathrm{geo}},
  \label{eq:ctrl-alpharef}
\end{equation}
turns \eqref{eq:ctrl-s1-torque} into
\begin{equation}
  \tau = \tau_{\mathrm{cur},f} + J\left(\alpha_{\mathrm{ref}} - \alpha_f\right),
  \label{eq:ctrl-indi-tau}
\end{equation}
the angular-acceleration increment of Tal and Karaman~\cite{tal2018indi}: the controller compares
the angular acceleration it wants with the one it gets, and changes the torque currently applied by
the difference. Both increments correct relative to what is actually being achieved, so they reject
a disturbance that appears in the measured accelerations without a model of it. Tal and Karaman
state the assumptions this needs: the disturbance changes slowly compared with the filters and the
sample period, and the measured quantities are filtered alike~\cite{tal2018indi}. The filters are
the source of the delay identified in Section~\ref{sec:model-measure-fail}.

\paragraph{Reference.} The structure is that of the INDI controller evaluated by Cobo-Briesewitz
et al.: the geometric law of Lee et al.\ with the residual force and moment, estimated online from
the inertial measurements and the rotor speeds, subtracted as feedforward terms, their eqs.~(8a),
(8b), (10a) and~(10b)~\cite{cobobriesewitz2026lindi}. The incremental principle and the requirement
of identical filters follow Tal and Karaman~\cite{tal2018indi} and, for the cascade of a
linear-acceleration loop over the attitude loop, Smeur et al.~\cite{smeur2017cascadedindi}.

\paragraph{Deviations.} From Cobo-Briesewitz et al.: the moment estimate omits the gyroscopic
term of their eq.~(10b), so the correction is $\tau_{\mathrm{cur},f} - J\alpha_f$ rather than the full
residual moment, and the gyroscopic coupling, which \eqref{eq:ctrl-tau-geo} already compensates,
enters the correction a second time; the vertical position integral and the attitude integral of
the common controller have no counterpart in their law; there is no payload term; and the rotor
speeds come from an optical rotor-speed sensor with one thrust constant for all rotors. From Tal and
Karaman: the rotor speeds are not measured by encoders on the motors~\cite{tal2018indi}, and Tal and
Karaman note that the incremental structure removes the need for integral action in their
design~\cite{tal2018indi}, whereas here the integrals of the common controller remain. No claim is
made that the flown law reproduces the tracking performance reported in either work.
% CHANGED: code, controller_omar_indi.c L408 (indi_moments = tau_rpm_f - tau_gyro_f; no omega x J omega
%   term), L131-136 (identical filters). Cobo-Briesewitz eq. (10b): tau_a = J omega_dot - J omega x omega
%   - tau_u (R07, PDF Sec. 4.2).

\paragraph{Requirements.} The rotor-speed source, without which switches A and B are inert; the
method needs no training data.

% -------------------------------------------------------------------------------------
\section{Strategy 2: geometric with a learned residual}
\label{sec:ctrl-s2}

% CHANGED (round 2): the learned residual on the common controller.
\paragraph{The law.} Switch~C. The deep-sets prediction $\hat{a}_{\mathrm{res}}$ of
\eqref{eq:model-deepsets} is subtracted in \eqref{eq:ctrl-ades}. The network is evaluated on the
vehicle at \SI{100}{\hertz} from the relative states of the neighbours: their positions arrive over
the radio with time stamps, and their velocities are obtained by finite differences
(Section~\ref{sec:model-deepsets}). The prediction is a vertical acceleration; the horizontal
residual is not compensated. The weights are trained offline on the dedicated training flights of
Strategy~0 and are fixed in the flight software for the whole comparison.
% CHANGED: code, ns2_residual.c (ns2_update every rnn.div = 10 main-loop ticks; ns2_eval peer positions
%   with time stamps from peer_get_all, finite differences with first-sample / repeated-time-stamp /
%   gap >= 500 ms rules; output z only, clamp 8 m/s^2); weights flash-resident (ns2_weights.c).
% CONFIRM-AFTER-LAB (c): the execution time of one evaluation on the vehicle is measured at the
%   hardware gate (rnn.us_last / us_max); no figure is stated before that.
Its execution time on the vehicle is measured at the hardware gate.

The prediction is computed and recorded \emph{whether switch~C is set or not}. On a flight in
which it does not drive the controller, predicted and measured residuals are two independent
quantities, and their comparison measures the model's accuracy. Once the prediction drives the
controller, it has itself moved the vehicle, and the comparison no longer isolates model quality.
Both kinds of flight are therefore needed, and Chapter~\ref{ch:setup} schedules them.

\paragraph{Training data.} The model is trained on dedicated flights under Strategy~0, never on the
comparison flights: a grid of vertical separations, lateral offsets and speeds in the scenarios of
the library, and randomised waypoint flights planned offline, in which the two vehicles are planned
together, the upper one stays at least \SI{0.25}{\metre} above the lower one, and the goals of the
upper vehicle are biased towards the lower vehicle's horizontal position so that the pair spends most
of the time inside the network's input gate. Every such flight passes the same safety check as the
scenarios of the library. The labels are the vertical component of the measured residual
\eqref{eq:model-ares}, computed after the flight from the recorded channels with the same filters as
switch~A, since under Strategy~0 the controller does not compute it.
% CONFIRM-AFTER-LAB (b): the number of training flights and every figure about the retrained model
%   (data volume, training result) are stated only after the training flights; the plan
%   (docs/plan_ns2_retrain_data_2026-10-14.md: 66 flights) is not a result and is not quoted.
% CHANGED: code/plan, crazyswarm2 scenarios.py R1 (offline-planned random waypoints, seeded,
%   safety.check, dz_min 0.25, stack bias); tools/residual/omar_labels.py (offline label with the
%   controller's filter chain).
% VERIFY: which rotor-speed source the offline labels use (optical sensor, as switch A, or the motor
%   controllers' telemetry) is open (round2_open_items.md).

\paragraph{Reference.} Predicting the interaction from relative states with a
permutation-invariant network and applying it as feedforward follows the Neural-Swarm
line~\cite{shi2020neuralswarm,shi2022neuralswarm2}. The closest published instance of the structure
used here, a learned residual added as feedforward to a geometric $SE(3)$ controller, is a
secondary configuration in the hardware experiments of Hsieh et al., whose main controller is the
feedback-linearising one of Strategy~4~\cite{hsieh2026flatness}.

\paragraph{Deviations.} The network architecture is Neural-Swarm2's
(Section~\ref{sec:model-deepsets}), but four things differ from that paper. First, Neural-Swarm2
applies spectral normalisation, which its abstract associates with stability and generalisation
guarantees and on which the bounded-error assumption of its analysis rests~\cite{shi2022neuralswarm2};
the model here is trained without it, and those guarantees are not claimed. Second, Neural-Swarm2
uses the learned residual in motion planning as well as control~\cite{shi2022neuralswarm2}, whereas
here the trajectories are planned in advance and the residual is used for control only. Third,
Neural-Swarm2's tracking controller has its own attitude law and an explicit compensation of motor
delay~\cite{shi2022neuralswarm2}; here the prediction enters the common controller unchanged.
Fourth, the model is trained on this project's two-vehicle data, it predicts only the vertical
component, its output is bounded at \SI{8}{\metre\per\second\squared}, at most three neighbours
are considered, and it is evaluated at \SI{100}{\hertz}. This strategy represents the predictive
family and should be read as such, not as an evaluation of Neural-Swarm2.

\paragraph{Requirements.} Neighbour relative states and training data. No rotor-speed sensing at
run time; the training labels require it on the vehicles that collect the data.

% -------------------------------------------------------------------------------------
\section{Strategy 3: INDI with a learned residual}
\label{sec:ctrl-s3}

% NEW (round 2): the hybrid is in the comparison (author decision 2).
\paragraph{The law.} Switches A, B and C together: the measured residual force and moment of
Strategy~1 and the learned residual of Strategy~2, on the common controller. In the position loop,
\eqref{eq:ctrl-ades} subtracts both the filtered measurement $a_{\mathrm{res},f}$ and the prediction
$\hat a_{\mathrm{res}}$; in the attitude loop, \eqref{eq:ctrl-s1-torque} applies.

\paragraph{Two estimates of one disturbance.} The measured and the predicted term act on the same
vertical disturbance, the residual force of the interaction, and the question is how they combine.
The measured residual is formed from the rotor speeds the vehicle actually has
(Section~\ref{sec:model-measure}), so it equals the residual force at the current state whatever
the controller commands: a prediction that raises the thrust raises $a_{\mathrm{rot}}$ and
$a_{\mathrm{meas}}$ alike and leaves $a_{\mathrm{res}}$ unchanged. The two terms therefore add.
The increment form shows what that means. With switch~C, the force increment of
Section~\ref{sec:ctrl-s1} becomes
\begin{equation}
  f R e_3 = f_{\mathrm{cur},f} R e_3 + m\left(a_{\mathrm{des}} - \hat a_{\mathrm{res}} - a_{\mathrm{meas},f}\right),
  \label{eq:ctrl-s3-increment}
\end{equation}
and when the thrust has settled, $f = f_{\mathrm{cur}}$, the filtered acceleration satisfies
$a_{\mathrm{meas},f} = a_{\mathrm{des}} - \hat a_{\mathrm{res}}$. For the INDI loop the prediction is
therefore not a second estimate of the disturbance, which the measurement has already removed, but
an offset of its acceleration reference by $-\hat a_{\mathrm{res}}$. In a quasi-static stack the
position loop answers that offset with a position error, and the vertical integral removes it at
steady state. Where the two estimates differ in time --- the measurement is delayed by its filters,
the prediction acts on the relative state as soon as it changes --- the prediction leads the
measurement, and that is where Strategy~3 can differ from Strategy~1. These are modelling
statements derived from \eqref{eq:model-ares} and \eqref{eq:ctrl-ades}, not results.
% DECISION D-12 (round 2, new): the author's instruction for this section said that "the measured
%   residual falls as the prediction takes effect, because the measurement is of the remaining
%   acceleration". With the residual formed from MEASURED rotor speeds (code: a_rpm from rpm.m1..m4,
%   controller_omar_indi.c L179-196, L250-262) this does not hold: compensation changes the thrust
%   but not the measured residual at a given state (Ch3 Sec. model-measure, fourth paragraph; Shi et
%   al. 2020 Sec. III-C). The derivation above follows the code. In NA-INDI, Cobo-Briesewitz et al.
%   split f_a = f_a,NN + f_a,INDI and let INDI "reason about the remaining mismatch" (R07, PDF
%   Sec. 4.4), i.e. the two terms are not simply added. Options and recommendation in
%   round2_open_items.md. The text above states the code's behaviour; it does not decide the design.

\paragraph{Reference.} Cobo-Briesewitz et al.\ propose combining a learned residual with INDI's
measured one on a single vehicle, writing the residual as the sum of a learned part and a part the
incremental term handles, so that the incremental term reasons about the remaining
mismatch~\cite{cobobriesewitz2026lindi}. Strategy~3 is a hybrid of the same two information
sources on a team, where the learned part is predicted from the neighbours' relative state.

\paragraph{Deviations.} The learned part is the deep-sets model of Strategy~2, predicted from the
neighbours' relative states and limited to the vertical component, not a model of the vehicle's own
dynamics; and the measured residual is not reduced by the prediction before it is used, so the two
terms add, with the consequence derived above. The deviations of Strategies~1 and~2 apply as well.

\paragraph{Requirements.} Those of Strategies~1 and~2 together: the rotor-speed source, the
neighbours' relative states and training data.

% -------------------------------------------------------------------------------------
\section{Strategy 4: feedback linearisation with a flatness-preserving residual}
\label{sec:ctrl-s4}

% CHANGED (round 2): renumbered from Strategy 3 to Strategy 4.
% DECISION D-05 (H9): implementation status. Kept neutral; nothing about the authors' code is stated.
\paragraph{Implementation status.} \textbf{Not implemented} in the flight software. The strategy is
described as a reference design. Its absence is a scope limitation, not a failure of the
comparison (the discussion chapter).

\paragraph{Reference.} Hsieh et al.~\cite{hsieh2026flatness}, building on the theory of Yang et
al.~\cite{yang2025flatnessresiduals}. The distinguishing property is where the prediction enters:
not as a feedforward term added to a desired acceleration, but inside a feedback-linearising law
whose validity depends on the augmented system remaining differentially flat. A generic learned
residual can destroy that flatness; residuals with a lower-triangular structure, a structured
correction inside the state derivative, not a free wrench, preserve it, and with it the original
flat output~\cite{yang2025flatnessresiduals}. That theory is established in simulation on a planar
quadrotor; the hardware demonstration is Hsieh et al.'s.
In Hsieh et al.\ the structure takes a simple form: the learned residual may depend only on the
positions and velocities of the vehicles, not on thrust or attitude, so that thrust and attitude
can still be computed algebraically from the flat outputs, position and yaw, and their
derivatives~\cite{hsieh2026flatness}.

Its relevance to this comparison is that it is the strongest published counter-argument to treating
``predictive compensation'' as one method. Hsieh et al.'s abstract reports a \SI{31}{\percent}
reduction in average tracking error. On hardware the body reports \SI{24}{\percent} RMSE for
ROM-FBL over nominal FBL and a further \SI{29}{\percent} for learned FBL over ROM-FBL, and an
average RMSE of about \SI{5}{\centi\metre} that the authors describe as matching the reported
performance of the nonlinear model predictive controller of Chee et al.; the direct comparison with
that controller, using the same learned residual, is in simulation, where the flatness-based
controller needs about one twentieth of the computation time~\cite{hsieh2026flatness}. The abstract
adds that tight formation flight was achieved with under \SI{30}{\second} of training data and a
\SI{5}{\milli\second} loop rate; on hardware that loop ran \emph{off-board}, on a laptop streaming
commands over the radio~\cite{hsieh2026flatness}.
Those are the figures a predictive strategy in this class should be measured against. None of them
was produced by a controller running on the vehicle.

\paragraph{Requirements.} Neighbour relative states; training data. No rotor-speed sensing at run
time; as for Strategy~2, the training labels require it.

% -------------------------------------------------------------------------------------
\section{What makes the comparison fair, and what it does not guarantee}
\label{sec:ctrl-fairness}

% CHANGED (round 2): the fairness argument rewritten for one controller (author decision 6).
\paragraph{What is shared.} The airframe, the trajectories, the formations and separations, the
estimator, the residual definition, and, because Strategies~0--3 are one controller, the position
gains, the vertical integral, the attitude gains and the attitude integral, the filters, and the
mass, inertia and thrust constant of Table~\ref{tab:ctrl-params}. The vertical integral is active in
all four, so that a constant height offset is removed in each of them and a difference in steady
height does not arise merely because one strategy lacks the integral.

\paragraph{What differs by design.} Only the switches. A difference in outcome between two of
Strategies~0--3 is therefore attributable to the terms the switches add: the measured residual force
and moment (A, B), the learned residual (C), or their combination.

\paragraph{What this does not guarantee.} Four limitations are stated before results are presented
rather than after. First, the shared gains are one set for all four strategies. A gain set suits
some switch combinations better than others; a strategy may be disadvantaged by gains that were not
chosen for it, and the comparison does not separate that effect from the effect of the switches.
Second, the strategies follow published \emph{families}, not papers reproduced exactly; the
deviations are listed per strategy, and the results are statements about these switch combinations
of one controller, not about the families in general. Third, the interaction of the measured and
the predicted term in Strategy~3 is the one derived in Section~\ref{sec:ctrl-s3}; a hybrid that
combines them differently could behave differently. Fourth, Strategy~4 is a different controller
and not implemented; if it were added, the shared-controller argument would not cover it, and
fairness towards it would rest on the experimental conditions alone.

% -------------------------------------------------------------------------------------
\section{Summary}
\label{sec:ctrl-summary}

% CHANGED (round 2).
Strategies~0--3 are one geometric controller with three switches: none (the uncompensated
baseline), the measured residual force and moment (INDI), the learned residual (geometric + learned
residual), and both (the hybrid). They share every gain, filter and physical parameter, and the
upper vehicle of every formation flies the baseline. Strategy~4, a feedback-linearising law with a
flatness-preserving learned residual, is a reference design and is not implemented.
Chapter~\ref{ch:setup} describes the experimental design; the two- and multi-robot results chapters
report what is flown.
